
\documentclass[a4paper,10pt]{report}

%\usepackage[landscape]{geometry}
\usepackage[cm-default]{fontspec}
\setromanfont{FreeSerif}
\setsansfont{FreeSans}
\setmonofont{FreeMono}
\usepackage[utf8x]{inputenc}
\usepackage{fontspec}
\usepackage{xunicode}
\usepackage{xltxtra}
\usepackage{xgreek}
\usepackage{amsmath}
\usepackage{unicode-math}
\usepackage{ulem}
\usepackage{color}
\usepackage{verbatim}
\usepackage{nopageno}
\usepackage{graphicx}
\usepackage{textpos}
\setlength{\TPHorizModule}{1cm}
\setlength{\TPVertModule}{1cm}
%\usepackage[colorgrid,texcoord]{eso-pic}
\usepackage[outline]{contour}
\usepackage{wrapfig}
\usepackage{url}
\usepackage{color}
\usepackage{tikz}
\usepackage{fancybox,fancyhdr}
\usepackage{subfigure}
\usepackage{pstricks}
\usepackage{epsfig}
\usepackage{multicol}
\usepackage{listings}
\usepackage{enumerate}
\usepackage{hyperref}
\hypersetup{
  bookmarks=true,
  bookmarksopen=true,
  pdfborder=false,
  pdfpagemode=UseNone,
  raiselinks=true,
  pdfhighlight={/P},
  colorlinks,
  citecolor=black,
  filecolor=black,
  linkcolor=black,
  urlcolor=black
}


\usepackage{fontspec}
\usepackage{xunicode}
\usepackage{xltxtra}


% Margins
%\setlength{\textwidth}{24cm}

%\setlength{\voffset}{0in}
%\setlength{\textheight}{6.8in}
% Colors
\definecolor{cbrown}{rgb}{0.49,0.24,0.07}
\definecolor{cmpez}{rgb}{0.92,0.88,0.79}
\definecolor{cmpez2}{rgb}{0.62,0.33,0.34}
\definecolor{cwhite}{rgb}{1,1,1}
\definecolor{cblack}{rgb}{0,0,0}
\definecolor{cred}{rgb}{0.9,0.15,0.15}
\definecolor{clblue}{rgb}{0.57,0.85,0.97}
\definecolor{corange}{rgb}{0.86,0.49,0.18}
\definecolor{clorange}{rgb}{0.95,0.84,0.61}
\definecolor{cyellow}{rgb}{1,0.95,0.16}
\definecolor{cgreen}{rgb}{0.69,0.8,0.31}
\definecolor{clbrown}{rgb}{0.78,0.67,0.43}
\definecolor{cmagenta}{rgb}{0.79,0.05,0.54}
\definecolor{cgray}{rgb}{0.76,0.73,0.63}
\definecolor{bluesite}{rgb}{0.05,0.43,0.67}
\definecolor{redsite}{rgb}{0.90,0.15,0.15}
\definecolor{backsite}{rgb}{0.07,0.07,0.07}
% ----------- from Costas -----------------------------

%---------------   Main Fonts    -----------------%   

% \setmainfont[Mapping=tex-text]{Calibri}
% \setmainfont[Mapping=tex-text]{Times New Roman}
%\setmainfont[Mapping=tex-text]{Droid Serif}
%\setmainfont[Mapping=tex-text]{Cambria}
%\setmainfont[Mapping=tex-text]{cm-unicode}
% \setmainfont[Mapping=tex-text]{Gentium}
% \setmainfont[Mapping=tex-text]{GFS Didot}
% \setmainfont[Mapping=tex-text]{Comic Sans MS}
 \setmainfont[Mapping=tex-text]{Ubuntu}
% \setmainfont[Mapping=tex-text]{Myriad Pro}
% \setmainfont[Mapping=tex-text]{CMU Concrete}
%\setmainfont[Mapping=tex-text]{DejaVu Sans}
%\setmainfont[Mapping=tex-text]{KerkisSans}
%\setmainfont[Mapping=tex-text]{KerkisCaligraphic}
% \setmainfont[Mapping=tex-text]{Segoe Print}
% \setmainfont[Mapping=tex-text]{Gabriola}

%------------------------------------------------%

% ------------ Mathematics Fonts ----------------%
\setmathfont{Asana-Math.ttf}
% -----------------------------------------------%
% Misc
\author{Κασωτάκη Ε. - Σμαραγδάκης Κ.}
\title{www.math24.gr}
% ---------------- Formation --------------------%


\setlength\topmargin{-0.7cm}
\addtolength\topmargin{-\headheight}
\addtolength\topmargin{-\headsep}
\setlength\textheight{26cm}
\setlength\oddsidemargin{-0.54cm}
\setlength\evensidemargin{-0.54cm}
%\setlength\marginparwidth{1.5in}
\setlength\textwidth{17cm}

\RequirePackage[avantgarde]{quotchap}
\renewcommand\chapterheadstartvskip{\vspace*{0\baselineskip}}
\RequirePackage[calcwidth]{titlesec}
\titleformat{\section}[hang]{\bfseries}
{\Large\thesection}{12pt}{\Large}[{\titlerule[0.9pt]}]
%--------------------------------------------------%
\usepackage{draftwatermark}
\SetWatermarkText{www.math24.gr}
\SetWatermarkLightness{0.9}
\SetWatermarkScale{0.6}
%--------------------------------------------------%


\usepackage{xcolor}
\usepackage{amsthm}
\usepackage{framed}
\usepackage{parskip}

\colorlet{shadecolor}{bluesite!20}

\newtheorem*{orismos}{Ορισμός}

\newenvironment{notation}
  {\begin{shaded}\begin{theorem}}
  {\end{theorem}\end{shaded}}




% Document begins
\begin{document}
\pagestyle{fancy}
\fancyhead{}
\fancyfoot{}
\renewcommand{\headrulewidth}{0pt}
\renewcommand{\footrulewidth}{0pt}


\fancyhead[LO,LE]{
 %\textblockcolor{backsite}
 %\begin{textblock}{5}(-2,-0.55)
  %\rule{0cm}{1cm}
 %\end{textblock}
 \textblockcolor{bluesite}
 \begin{textblock}{5}(-1.5,-0.55)
  \rule{0cm}{1cm}
 \end{textblock}
 %\textblockcolor{bluesite}
 %\begin{textblock}{14}(5.5,-0.55)
 % \rule{0cm}{1cm}
 %\end{textblock}
 \begin{textblock}{0}(-1,-0.25)
 \color{cwhite} \begin{Large}www.math24.gr\end{Large}
 \end{textblock}
\begin{textblock}{0}(16,-0.4)
 \color{cwhite} \includegraphics[height=1.5cm]{math24_logo.png}
 \end{textblock}
\textblockcolor{white}
\begin{textblock}{17}(-1,28)
 \color{backsite} \begin{small}Copyright \textcopyright 2012, Κασωτάκη Ε.(ikasotaki@gmail.com) - Σμαραγδάκης Κ.(kesmarag@gmail.com)\end{small}
 \end{textblock}
}
 
\begin{shaded}
\begin{center}
\huge \textbf{Φυλλάδιο Ασκήσεων}\\
%Πρόσθεση ρητών αριθμών
\end{center} 
\textbf{Μαθηματικά Α' Γυμνασίου} \hfill \textbf{Ημερομηνία Παράδοσης : \hspace{2em} }
\subsection*{Ονοματεπώνυμο :\hfill  \hspace{5em}}
\end{shaded}
\vspace{2em}
\begin{itemize}
 \item Ορισμός αριθμητικής παράστασης
 \item Προτεραιότητα των πράξεων
\end{itemize}

\section*{Θεωρία - Αριθμητική Παράσταση \hfill \small{}}
\textbf{Αριθμητική παράσταση} λέγεται κάθε σειρά αριθμών που συνδέονται μεταξύ τους με τα σύμβολα των πράξεων.\\
\textbf{Π.χ} $2+7$ είναι αριθμητική παράσταση\\
\textbf{Π.χ} $3\cdot7+4^{2}-2$ είναι αριθμητική παράσταση


\section*{Θεωρία - Προτεραιότητα των Πράξεων  \hfill \small{}}
\begin{itemize}
 \item Σε μία αριθμητική παράσταση εκτελούμε τις πράξεις σύμφωνα με την παρακάτω σειρά:
       \begin{itemize}
        \item \textbf{1\textdegree  βήμα:} Υπολογίζουμε τις δυνάμεις
        \item \textbf{2\textdegree  βήμα:} Κάνουμε τους πολλαπλασιασμούς και τις διαιρέσεις
        \item \textbf{3\textdegree  βήμα:} Κάνουμε τις προσθέσεις και τις αφαιρέσεις
       \end{itemize}
 \item Αν στην αριθμητική παράσταση υπάρχουν παρενθέσεις κάνουμε πρώτα τις πράξεις με την παραπάνω σειρά
 \item Αν κάποιο από τα παραπάνω βήματα "δεν υπάρχει" προχωράμε στο επόμενο.\\
       \textbf{Π.χ } Στην παράσταση $7\cdot4+8$ δεν υπάρχουν δυνάμεις άρα το 1\textdegree  βήμα παραλείπεται 
        και προχωράμε στο επόμενο.
\end{itemize}
\textbf{Παράδειγμα}\\
Θέλουμε να υπολογίσουμε την τιμή της αλγεβρικής παράστασης: $Α=2+5\cdot 1^{10}$
 \begin{itemize}
        \item \textbf{1\textdegree  βήμα:} Πρώτα υπολογίζουμε το $1^{10}$ και βρίσκουμε ότι$1^{10}=1$.\\
               Άρα $Α=2+3\cdot5\cdot 1^{10}=2+3\cdot5\cdot1$
        \item \textbf{2\textdegree  βήμα:} Στη συνέχεια υπολογίζουμε το γινόμενο $5\cdot1$ και βρίσκουμε ότι 
         $5\cdot1=5$.\\
          Άρα $Α=2+5\cdot1=2+5$
        \item \textbf{3\textdegree  βήμα:} Τέλος υπολογίζουμε το άθροισμα $2+5$ και βρίσκουμε ότι $2+5=7$.\\
         Άρα $Α=7$
\end{itemize}


\section*{Άσκηση 1  \hfill \small{}25  μονάδες}
Να αντιστοιχίσετε κάθε στοιχείο της αριστερής στήλης με ένα στοιχείο της δεξιάς στήλης
\begin{itemize}
\begin{multicols}{2}
 \item $2+3\cdot4$
 \item $(2+3)\cdot4$
 \item $2\cdot(3+4)-2$
 \item $2\cdot (3+4-2)$
 \item $(2+3)-(4-2)$
 \item $3$
 \item $10$
 \item $12$
 \item $14$
 \item $20$
\end{multicols}
\end{itemize}


\section*{Άσκηση 2  \hfill \small{}25  μονάδες}
Να αντιστοιχίσετε κάθε στοιχείο της αριστερής στήλης με ένα στοιχείο της δεξιάς στήλης
\begin{itemize}
\begin{multicols}{2}
 \item $2^{2}\cdot2-4$
 \item $2^{2}\cdot(4-2)$
 \item $2^{2}-4$
 \item $2^{2}\cdot4\cdot2$
 \item $2^{2}+4+2$
 \item $0$
 \item $32$
 \item $10$
 \item $8$
 \item $4$
\end{multicols}
\end{itemize}

\section*{Άσκηση 3  \hfill \small{50 μονάδες}}
Να υπολογίσετε τις τιμές των παρακάτω αριθμητικών παραστάσεων:
\begin{enumerate}[1)]
 \item $7\cdot2-6$
 \item $6^{2}-2\cdot6$
 \item $3\cdot3^{2}-2^{2}\cdot5$
 \item $10^{2}-2\cdot5\cdot3\cdot1^{7}$
 \item $3-2+5-2^{2}$ 
 \item $(5-2)^{2}-2\cdot3$
 \item $10\cdot(4^{2}-2\cdot3)-(2\cdot5)^{2}$
 \item $8^{2}+6^{2}-2\cdot5\cdot5+1^{8}\cdot0$
 \item $2\cdot(2^{2}\cdot3-2\cdot3)^{2}$
 \item $2^{2}+2^{2}\cdot2^{2}-2^{2}$
\end{enumerate}



 








%\let\thefootnote\relax\footnotetext{Επιμέλεια : Κασωτάκη Ειρήνη (ikasotaki@gmail.com),    Σμαραγδάκης Κώστας (kesmarag@gmail.com)} 


\end{document}
